\chapter{Setup and troubleshooting}

\section{Required tools}

The runnable code requires Python 3.12 or newer. The PDF build additionally
requires XeLaTeX, \code{latexmk}, and the Noto Serif, Noto Sans, and JetBrains
Mono fonts. No third-party Python package is required.

Confirm Python:

\begin{wiremessage}
python --version
\end{wiremessage}

On systems where \code{python} refers to an older interpreter, try
\code{python3}. On Windows, the Python Launcher commonly uses \code{py -3.12} or
newer. Use the same launcher consistently in later commands.

\section{Optional virtual environment}

There are no packages to install, but a virtual environment can keep tooling
habits consistent.

\begin{center}
\begin{tabularx}{\textwidth}{@{}>{\bfseries}p{32mm} X X@{}}
\toprule
Task & Linux/macOS shell & Windows PowerShell \\
\midrule
Create & \code{python -m venv .venv} & \code{py -m venv .venv} \\
Activate & \code{source .venv/bin/activate} & \path{.venv\Scripts\Activate.ps1} \\
Deactivate & \code{deactivate} & \code{deactivate} \\
\bottomrule
\end{tabularx}
\end{center}

If PowerShell execution policy blocks activation, the environment can still be
used by invoking \path{.venv\Scripts\python.exe} directly. Do not weaken a
machine-wide policy merely for this project.

\section{Repository map}

\begin{itemize}
  \item \filepath{book}: LaTeX source and generated build intermediates;
  \item \filepath{code/stages}: nine frozen learning milestones and snapshot tests;
  \item \filepath{code/final/tinyhttp}: reusable educational framework;
  \item \filepath{code/final/task\_api}: SQLite application and CLI;
  \item \filepath{code/final/tests}: parser, application, API, and socket tests;
  \item \filepath{dist}: approved compiled PDF;
  \item \filepath{build.py}: cross-platform test/build helper.
\end{itemize}

\section{Run tests}

From the repository root:

\begin{wiremessage}
python build.py test
\end{wiremessage}

To isolate suites:

\begin{wiremessage}
cd code/stages
python -m unittest discover -s tests -v

cd ../final
python -m unittest discover -s tests -v
\end{wiremessage}

Local socket tests may be restricted in a sandboxed development environment.
Grant only loopback-test capability; do not disable unrelated security controls.

\section{Build the book}

\begin{wiremessage}
python build.py book
\end{wiremessage}

The helper runs \code{latexmk -xelatex}, inspects the log for undefined references
and overfull boxes, and copies the final artifact to
\filepath{dist/build-http-from-scratch.pdf}. Use \code{python build.py all} for
tests plus PDF. Use \code{python build.py clean} to remove generated LaTeX files
and the distribution PDF.

The manuscript uses \code{listings}, not \code{minted}; shell escape and Pygments
are unnecessary. Code listings are read directly from the runnable source, so
the relative directory structure must remain intact.

\section{Run the API}

\begin{wiremessage}
cd code/final
python -m task_api --database tasks.db
\end{wiremessage}

Useful options:

\begin{wiremessage}
python -m task_api --host 127.0.0.1 --port 9000 \
  --database practice.db --workers 4 --no-keep-alive
\end{wiremessage}

The backslash line continuation is for a Unix-like shell. In PowerShell, put the
command on one line or use its backtick continuation carefully. The options are
otherwise identical.

Stop with Ctrl+C. The database file and possible \code{-wal}/\code{-shm} companion
files belong to the running SQLite database. Stop the process before moving or
deleting a practice database.

\section{Common failures}

\subsection{Address already in use}

Another process is listening on the selected port, or a previous process is
still running. Choose another port or stop the exact process you started. Do not
kill unrelated processes by a broad name match.

\subsection{Connection refused}

Confirm the server printed its listening URL, the client uses the same address
family and port, and the process has not exited with a traceback. A server bound
to loopback cannot be reached through another machine's network address.

\subsection{Client waits forever}

Inspect response framing. A wrong Content-Length can make a client wait for bytes
that will never arrive. In learning snapshots without timeouts, a server may also
wait for an incomplete request delimiter.

\subsection{400 invalid Host}

Raw HTTP/1.1 requests must contain exactly one nonempty Host field. Standard
clients add it automatically. When using a raw socket, include
\code{Host: localhost} followed by CRLF.

\subsection{411 on POST, PUT, or PATCH}

Body-required Task endpoints require Content-Length. Standard HTTP clients
usually calculate it when given bytes. A raw client must calculate the UTF-8 byte
length, not character count.

\subsection{415 Unsupported Media Type}

Set \code{Content-Type: application/json}. An \code{Accept} field describes the
desired response and does not describe the request body.

\subsection{SQLite database is locked}

Stop stray API processes using the file. Keep transactions short. The repository
configures a busy timeout, but a long-held external transaction can exceed it.
Do not edit the database with a tool that leaves a write transaction open.

\subsection{Font or LaTeX error}

Confirm XeLaTeX can locate all three configured fonts and that the TeX distribution
contains memoir, fontspec, TikZ, tcolorbox, listings, hyperref, bookmark, and
cleveref. The local build intentionally avoids automatic network installation.

\section{A debugging worksheet}

For any failure, record:

\begin{enumerate}
  \item exact command and current directory;
  \item Python and tool versions;
  \item selected hostname/address/port;
  \item raw request start line, fields, and declared body length;
  \item first observed error, status, or traceback;
  \item earliest layer that could produce that evidence;
  \item smallest reproducible case and one falsifiable next hypothesis.
\end{enumerate}

This worksheet prevents random edits and preserves evidence for mentor review.
